\section{Proofs for the Upper Bounds}
\label{app:complexity}

We prove \cref{lem:poly-tests}, \cref{prop:summary-search}, and
\cref{thm:complexity}.

\subsection{Polynomial-Time Tests}

\begin{proof}[Proof of \cref{lem:poly-tests}]
The table of \cref{sec:signed} gives $\psi$ with $2k$ variables and two
literals for each literal of $\phi$; the order dual reverses its inequalities
and exchanges its constants. The
closure is computed by adding reflexive, transitive, and decomposition
consequences \eqref{eq:g21} until none is new. At most $(2k)^2$ pairs are
added, and each round inspects polynomially many pairs of literals and
closure pairs. By \eqref{eq:g22} with an empty remainder, $\Up(i,u,v)$ holds
iff $D_\psi(u,z)$, $z\le f(\vec b)\in\psi$, and $D_\psi(b_i,v)$ for some
$z,\vec b$, and dually for $\Lo$. Hence the updates $E_i,G_i$ of
\eqref{eq:g37}, the tests of \eqref{eq:g34}, the initial state and transitions
\eqref{eq:g38}, rejection \eqref{eq:g40}, and admission \eqref{eq:g41} are
operations on subsets of the $2k$ signed variables, computed from the
closure in polynomial time. A label clash is a lower and an upper literal
with incompatible root labels whose roots are related by the closure
(\cref{sec:least}). A cycle clash \eqref{eq:g55} is a directed cycle in the
graph on the at most $(2k)^2$ closure pairs used in the proof of
\cref{thm:finite-sat}; its edges need only single-letter path judgments, so
cyclicity is decided in polynomial time. Finally, $\Sat_c(\phi)$ and
$\Sat_{\mathrm{reg},c}(\phi)$ hold iff $\psi$ has no label clash
(\eqref{eq:g19}, \cref{thm:sat}, and \eqref{eq:sat-reg}), and
$\Sat_{\mathrm{fin},c}(\phi)$ holds iff $\psi$ has neither a label clash nor a
cycle clash (\cref{lem:finite-signed,thm:finite-sat}).
\end{proof}

\subsection{Summaries Determine Derived Monitors}

The proof of \cref{prop:summary-search} uses the lettered monitors and
derivatives of \cref{sec:search}, and the summaries, column words, and
agreement of \cref{app:search}. Throughout, $\mathcal M$ is a lettered monitor
with state set $\mathcal Q$ and start state $\zeta_0$, $\ell$ is one of its
letters, and $[h]=\{1,\ldots,h\}$.

The summary of blocks $X_1,\ldots,X_h$ avoiding $\ell$, with respect to
$\mathcal M$ and $\ell$, is the tuple
$\Theta=(\mathit{exit},\mathit{term},\mathit{eq},\mathit{adm})$ of
\cref{app:search}, which the column automaton described there computes along
the column word of the tuple. A column whose entries are all absent changes no
component of a summary. Every valid column word whose tracks spell
$X_1,\ldots,X_h$ is the column word of the tuple preceded by such columns, so
the automaton computes $\Theta$ along every such word.

\begin{lemma}[Summary monitor]\label{lem:summary-monitor}
Let $\Theta=(\mathit{exit},\mathit{term},\mathit{eq},\mathit{adm})$ be the
summary of pairwise distinct blocks $X_1,\ldots,X_h$ avoiding $\ell$. Let
$\mathcal M[\ell,\Theta]$ be the lettered monitor over $[h]$ with state set
$\mathcal Q\sqcup\{\star\}$, start state $\star$, and, for $a,b\in[h]$ and
$\zeta,\zeta'\in\mathcal Q\sqcup\{\star\}$,
\begin{equation}\label{eq:summary-monitor}
\begin{aligned}
 \delta(\zeta,a)&=\mathit{exit}_a(\zeta),\\
 \Tfin(\zeta,a)&\iff\Tfin_{\mathcal M}(\mathit{exit}_a(\zeta),\ell)\land
   \neg\,\mathit{term}_a(\zeta,\mathit{exit}_a(\zeta)),\\
 \Lambda((\zeta,a),(\zeta',b))&\iff\mathit{adm}(a,b,\zeta,\zeta'),\\
 \Lfin((\zeta,a),(\zeta',b))&\iff\mathit{adm}(a,b,\zeta,\zeta')\lor
   \mathit{term}_a(\zeta,\mathit{exit}_b(\zeta')),
\end{aligned}
\end{equation}
where $\Tfin_{\mathcal M}$ is the target of $\mathcal M$; every relation
involving the label $\triangleright$ is false. Then
\[
 \mathcal M[\ell,\Theta]\equiv_{[h],X}\mathcal M_\ell,
\]
where $X$ denotes the map $a\mapsto X_a$.
In particular, $\Theta$ and the targets $\Tfin_{\mathcal M}(\cdot,\ell)$
determine the restriction of the derivative $\mathcal M_\ell$ to the letters
$X_1,\ldots,X_h$, up to the names of these letters.
\end{lemma}

\begin{proof}
The right-hand sides are defined: an exit state is a state of $\mathcal Q$,
also from the entry state $\star$, and $\mathit{term}_a(\zeta,\zeta'')$ is
recorded for every $\zeta''\in\mathcal Q$. Both monitors have the state set
$\mathcal Q\sqcup\{\star\}$ and the start state $\star$, and $X$ is
injective because the blocks are distinct. We compare
\eqref{eq:summary-monitor} with \cref{def:derivative} at the letters
$X_a,X_b$ and the states $\zeta,\zeta'$. The derived transition
$\delta_\ell(\zeta,X_a)$ is the exit state of the shifted block of $X_a$ read
from $\zeta$, that is, $\mathit{exit}_a(\zeta)$. The bit
$\mathit{term}_a(\zeta,\zeta'')$ states that
$\Lfin(\mathit{lab}_p(\zeta,X_a),(\zeta'',\ell))$ holds for some
$p\le|X_a|$. For $\zeta''=\delta_\ell(\zeta,X_a)$ this is the condition
excluded by $\Tfin_\ell(\zeta,X_a)$, whose other condition is
$\Tfin_{\mathcal M}(\delta_\ell(\zeta,X_a),\ell)$; this proves the second
line. For $\zeta''=\delta_\ell(\zeta',X_b)$ it is the second disjunct of
$\Lfin_\ell((\zeta,X_a),(\zeta',X_b))$. By definition,
$\mathit{adm}(a,b,\zeta,\zeta')$ states
$\Lambda_\ell((\zeta,X_a),(\zeta',X_b))$, which is also the first disjunct
of $\Lfin_\ell$. This proves the third and fourth lines. Finally, every
relation involving the start label $\triangleright$ of $\mathcal M_\ell$ is
false by \cref{def:derivative}. These are all the conditions of
\eqref{eq:monitor-agreement} for the letters of $[h]$ and the labels over
them.
\end{proof}

\begin{lemma}[Summaries under agreement]\label{lem:summary-agreement}
Let $\mathcal U\equiv_{E,\rho}\mathcal V$ with $\ell\in E$, and let
$X_1,\ldots,X_h$ be blocks over $E$ avoiding $\ell$. Then the summary of
$X_1,\ldots,X_h$ with respect to $\mathcal U$ and $\ell$ equals the summary
of $\rho^*(X_1),\ldots,\rho^*(X_h)$ with respect to $\mathcal V$ and
$\rho(\ell)$.
\end{lemma}

\begin{proof}
The image blocks avoid $\rho(\ell)$, since $\rho$ is injective on $E$, and
the column word of the image tuple is the letterwise image of the column
word of $X_1,\ldots,X_h$, with absent entries kept absent. We show by
induction on the columns that the two column automata, for $\mathcal U,\ell$
and for $\mathcal V,\rho(\ell)$, reach equal summaries. Both monitors have
the same states and the same start state, and summaries take values in
states and bits only. In the initial summary,
$\mathit{exit}_a(\zeta)=\delta(\zeta,\ell)$ agrees by the first line of
\eqref{eq:monitor-agreement}, $\mathit{exit}_a(\star)$ is the common start
state, and the $\mathit{term}$ and $\mathit{adm}$ bits evaluate $\Lfin$ and
$\Lambda$ on the labels $\mathit{lab}_0(\zeta)$ and $(\zeta'',\ell)$. The map
$\widehat\rho$ sends these labels to the corresponding labels for the marker
$\rho(\ell)$, so the two lines of \eqref{eq:monitor-agreement} for the
relations give equal bits. An update evaluates $\delta$ at a letter of $E$,
$\Lambda$ and $\Lfin$ on labels whose letters lie in $E$, and the equality of
two entries of the column. Agreement preserves the first two, and
injectivity of $\rho$ on $E$ preserves equality and inequality of letters.
Hence the updated summaries are equal.
\end{proof}

\subsection{The Search}

We restate the search of \cref{fig:summary-search} with the details that the
space bound needs. Fix $\mathcal M_0$, with state set $Q$ of size $N\ge1$,
start state $\zeta_0$, and finite alphabet $\Gamma$, and fix $d\ge0$. Put
\[
 g=\DictBound(N,d),\qquad \bar N=N+d+1,\qquad \bar g=g+|\Gamma|;
\]
thus $\bar N$ bounds the number of states, and $\bar g$ the number of
letters, of every monitor that the search stores.
The search keeps a level $j\le d$ and a \emph{current monitor}
$\mathcal W$: a lettered monitor with state set
$Q\sqcup\{\star_1,\ldots,\star_j\}$, of size $N+j$, over a \emph{current
alphabet} $\Gamma_{\mathcal W}$, which is $\Gamma$ for $j=0$ and $[h]$ for some
$h\le g$ for $j\ge1$. Initially $j=0$ and $\mathcal W=\mathcal M_0$. The
search repeats one of two moves.
\begin{itemize}[leftmargin=*]
\item \emph{Accept.} Choose $z\in\Gamma_{\mathcal W}$ and $1\le m'\le N+j$, and
  accept if $z^{m'}$ is a hole of $\mathcal W$. Every pair of cuts of a
  unary word is a comparison, so the test checks the target at the last
  label, $\Lfin$ between the label of every cut $s<m'$ and the last label,
  and $\Lambda$ between the labels of every two cuts $s<e<m'$
  (\cref{app:search}, paragraph on unary words). The labels are computed
  along the run of $\mathcal W$ on $z^{m'}$.
\item \emph{Derive}, if $j<d$. Choose a marker $\ell\in\Gamma_{\mathcal W}$ and a
  number $1\le h\le g$. Guess a column word with $h$ tracks over
  $\Gamma_{\mathcal W}\setminus\{\ell\}$, one column at a time. Keep one flag per
  track recording whether the track has started, and refuse a column with an
  absent entry on a started track, so that the word is valid. Count the
  columns with a binary counter and stop before $\mathit{col}(N+j,h)$
  columns. Maintain the summary of the tracks read so far by the column
  automaton of \cref{app:search} for $\mathcal W$ and $\ell$, starting from
  the summary of $h$ empty blocks. At the end, with final summary $\Theta$,
  require $\mathit{eq}(a,b)$ to be false for all $a\ne b$. Then replace
  $\mathcal W$ by $\mathcal W[\ell,\Theta]$ of \cref{lem:summary-monitor},
  naming its new state $\star_{j+1}$, and replace $j$ by $j+1$.
\end{itemize}
The tracks of the guessed column word spell blocks $Y_1,\ldots,Y_h$ over
$\Gamma_{\mathcal W}$ that avoid $\ell$. Their summary is the final summary
$\Theta$, and the final test makes them pairwise distinct.

\paragraph{Space}
The current monitor has $N+j<\bar N$ states and at most $\bar g$ letters.
Write each state with $\lceil\log_2\bar N\rceil\le\bar N$ bits and each
letter, or absent entry, with $\lceil\log_2(\bar g+1)\rceil\le\bar g$ bits.
The search stores the following items.
\begin{itemize}[leftmargin=*]
\item The current monitor: its transitions, targets, and the relations
  $\Lambda$ and $\Lfin$ on its at most $1+\bar N\bar g$ labels, in at most
  $\bar N^2\bar g+\bar N\bar g+2(1+\bar N\bar g)^2\le10\bar N^2\bar g^2$ bits.
\item During a derive move, the next monitor $\mathcal W[\ell,\Theta]$, with
  $N+j+1\le\bar N$ states and $h\le\bar g$ letters, in at most
  $10\bar N^2\bar g^2$ bits.
\item One summary, for $h\le\bar g$ blocks, $N+j+1\le\bar N$ entry states, and
  $N+j<\bar N$ states: $\mathit{exit}$ uses at most $h\bar N\cdot\bar N$ bits,
  $\mathit{term}$ at most $h\bar N^2$, $\mathit{eq}$ at most $h^2$, and
  $\mathit{adm}$ at most $h^2\bar N^2$, in total at most $4\bar N^2\bar g^2$
  bits.
\item The column counter. By \eqref{eq:col}, which is nondecreasing in both
  arguments, it counts below $\mathit{col}(\bar N,g)$, and
  \[
   \log_2\mathit{col}(\bar N,h)=(\bar N+1)h\log_2\bar N+\bar N(\bar N+1)h+h^2
   +(\bar N+1)^2h^2\le4(\bar N+1)^2h^2,
  \]
  using $\log_2\bar N\le\bar N$. The counter therefore has at most
  $1+16\bar N^2g^2\le17\bar N^2\bar g^2$ bits.
\item The current column and the track flags, in at most
  $h(\bar g+1)\le2\bar g^2$ bits.
\item The level $j$, the choices $\ell$, $h$, $z$, $m'$, and, for the accept
  test, two cuts, the states at three cuts, and a running state with its
  counter, recomputed from the start state: at most
  $9\bar N+3\bar g\le12\bar N^2\bar g^2$ bits.
\end{itemize}
These items take at most $55\bar N^2\bar g^2$ bits. The summary is updated
in place, column by column. Of the old summary, an update of a field
$\mathit{term}$, $\mathit{eq}$, or $\mathit{adm}$ reads only its own old bit
and the old exit states, and an update of a field $\mathit{exit}$ reads only
its own old value. So the search first updates the fields $\mathit{term}$,
$\mathit{eq}$, and $\mathit{adm}$, retaining the old exit states, and then
the fields $\mathit{exit}$. Alternatively, a second summary buffer of the
same size holds the new values. The remaining work space consists of the
position of the input head, the counters and indices that address entries of
the stored tables, the binary arithmetic that computes
$\mathit{col}(N+j,h)$ and compares the column counter with it, and the
evaluation of single entries of the update formulas and of
\eqref{eq:summary-monitor}. Every number involved is below
$\mathit{col}(\bar N,g)$ or below the size of the input or of a stored table,
so this work space takes $O(\bar N^2\bar g^2)$ bits. Hence the search uses at
most $\mathsf{Sp}=c_{\mathsf{Sp}}\bar N^2\bar g^2$ bits, for an absolute
constant $c_{\mathsf{Sp}}$; this is polynomial in $\bar N$, $g$, and
$|\Gamma|$. All bounds above use only the column bound \eqref{eq:col}, the
dictionary bound $g$, and the unary bound $m'\le N+j$.

\paragraph{Soundness}
We show that whenever the search is at level $j$ with current monitor
$\mathcal W$, there are markers
$\widetilde{\boldsymbol\ell}=(\widetilde\ell_0,\ldots,\widetilde\ell_{j-1})$
and a map $\rho_j$ from $\Gamma_{\mathcal W}$ to letters of level $j$,
injective with an image family valid for $\widetilde{\boldsymbol\ell}$, such
that
\[
 \mathcal W\equiv_{\Gamma_{\mathcal W},\rho_j}
 \mathcal M^{\widetilde{\boldsymbol\ell}}_j ,
\]
in the notation of \cref{app:search} for towers. Initially $\rho_0$ is the
identity on $\Gamma$ and the marker sequence is empty.

For a derive move, let $Y_1,\ldots,Y_h$ be the spelled blocks and $\Theta$
their summary. By \cref{lem:summary-monitor},
$\mathcal W[\ell,\Theta]\equiv_{[h],Y}\mathcal W_\ell$, where $Y$ is the map
$a\mapsto Y_a$. By \cref{lem:agreement-derivative}, applied to the family
$\{Y_1,\ldots,Y_h\}$ of blocks over $\Gamma_{\mathcal W}$ avoiding
$\ell\in\Gamma_{\mathcal W}$,
\[
 \mathcal W_\ell\equiv_{\{Y_a\},\rho_j^*}
 \bigl(\mathcal M^{\widetilde{\boldsymbol\ell}}_j\bigr)_{\rho_j(\ell)},
\]
and the image blocks avoid $\rho_j(\ell)$. Put
$\widetilde\ell_j=\rho_j(\ell)$ and $\rho_{j+1}(a)=\rho_j^*(Y_a)$. Composing
the two agreements gives the invariant at level $j+1$; appending a marker
does not change the lower monitors of the tower. The map $\rho_{j+1}$ is
injective because the blocks are distinct and $\rho_j^*$ is injective on
words over $\Gamma_{\mathcal W}$. Each image block avoids $\widetilde\ell_j$
and consists of valid level-$j$ letters, so the image family is valid.

If the search accepts with $z^{m'}$, agreement preserves holes of words over
$\Gamma_{\mathcal W}$, so $\rho_j(z)^{m'}$ is a hole of
$\mathcal M^{\widetilde{\boldsymbol\ell}}_j$. Expand it along
$\widetilde\ell_{j-1},\ldots,\widetilde\ell_0$. At each level the word is a
nonempty word of blocks avoiding that level's marker, by validity, so
\cref{lem:der-holes} makes its expansion a hole of the monitor one level
below. This yields a nonempty hole $w$ of $\mathcal M_0$. Each expansion ends
with its marker, and its blocks avoid that marker. As in the proof of
\cref{thm:search}, cutting $w$ at its successive markers recovers the
expanded words, so the canonical hierarchy of $w$ reaches the unary word
$\rho_j(z)^{m'}$ at level $j$, and $\depth(w)\le j\le d$.

\paragraph{Completeness}
Suppose that $\mathcal M_0$ has a nonempty hole of depth at most $d$. By
\cref{lem:dictionary-bound}, it has a nonempty hole $w$ of some depth
$\Delta\le d$ whose dictionaries satisfy $g_i\le g$ for $1\le i\le\Delta$. Let $u_i$,
$\ell_i$, $\mathcal A_i$, and $\mathcal M_i$ be its hierarchy, markers,
dictionaries, and derived monitors, so that $\mathcal M_{i+1}=(\mathcal
M_i)_{\ell_i}$ and $u_i$ is a hole of $\mathcal M_i$ (\cref{lem:der-holes}).
We follow a run that reaches, at each level $j\le\Delta$, a current monitor
$\mathcal W_j$ and a map $\pi_j$ with
\[
 \mathcal W_j\equiv_{\Gamma_{\mathcal W_j},\pi_j}\mathcal M_j,\qquad
 \mathcal A_j\subseteq\pi_j(\Gamma_{\mathcal W_j}).
\]
At $j=0$ take $\mathcal W_0=\mathcal M_0$ and $\pi_0$ the identity on
$\Gamma$.

Let $j<\Delta$, so a derive move is allowed. The $h=g_{j+1}\le g$ letters
$X_1,\ldots,X_h$ of $\mathcal A_{j+1}$ are distinct blocks of $u_j$; their
letters lie in $\mathcal A_j$, and they avoid $\ell_j\in\mathcal A_j$.
Choose $\ell=\pi_j^{-1}(\ell_j)$ and $Y'_a=(\pi_j^{-1})^*(X_a)$, which avoid
$\ell$. By \cref{lem:summary-agreement}, the summary $\Theta$ of
$Y'_1,\ldots,Y'_h$ with respect to $\mathcal W_j$ and $\ell$ equals that of
$X_1,\ldots,X_h$ with respect to $\mathcal M_j$ and $\ell_j$; in particular,
its equality fields declare the blocks pairwise distinct. The current
monitor has $N+j$ states, so by column shortening (\cref{app:search},
paragraph on shortening), some valid column word with fewer than
$\mathit{col}(N+j,h)$ columns has the same final summary $\Theta$. The
search chooses $\ell$ and $h$, guesses this column word, and passes the final
test. Its new monitor $\mathcal W_{j+1}=\mathcal W_j[\ell,\Theta]$ depends
only on $\Theta$. Since $\Theta$ is also the summary of the distinct blocks
$Y'_a$, \cref{lem:summary-monitor} gives
$\mathcal W_{j+1}\equiv_{[h],Y'}(\mathcal W_j)_\ell$, where $Y'$ is the map
$a\mapsto Y'_a$. \Cref{lem:agreement-derivative} gives
$(\mathcal W_j)_\ell\equiv_{\{Y'_a\},\pi_j^*}(\mathcal M_j)_{\ell_j}
=\mathcal M_{j+1}$. Composing, $\mathcal W_{j+1}\equiv_{[h],\pi_{j+1}}\mathcal
M_{j+1}$ for $\pi_{j+1}(a)=\pi_j^*(Y'_a)=X_a$, a bijection onto
$\mathcal A_{j+1}$.

At level $\Delta$, the word $u_\Delta=z^L$ is a unary hole of $\mathcal M_\Delta$,
which has $N+\Delta$ states. By unary shortening (\cref{app:search}, paragraph on
unary words), $z^{m'}$ is a hole of $\mathcal M_\Delta$ for some
$1\le m'\le N+\Delta$. Write $z=\pi_\Delta(z')$. By agreement, $z'^{m'}$ is a hole of
$\mathcal W_\Delta$, and the search accepts. For $\Delta=0$ this is the accept move at
level zero.

\paragraph{Deterministic Space}
The configurations of the search consist of the items counted above, a
control state, and the input head position. The bound
$\mathsf{Sp}=c_{\mathsf{Sp}}\bar N^2\bar g^2$ is at least the logarithm of
the input size, and it can be computed from $N$, $d$, and $|\Gamma|$ in space
$O(\mathsf{Sp})$. Savitch's theorem~\citep{Savitch70} simulates a
nondeterministic machine with a space bound of this kind, at least
logarithmic and constructible, by a deterministic machine in the square of
that space. Hence the search runs in deterministic space
$O(\mathsf{Sp}^2)=O(\bar N^4\bar g^4)$. Since $g\ge1$, we have
$\bar g\le g(1+|\Gamma|)$, and $g\le2^{\Xi(N,d)}$ by
\cref{lem:elementary-estimate}. Moreover,
\[
 \Xi(N,d)=3N^2+4N(N+d)(d+1)+3N+3d+4\le3\bar N^2+4\bar N^3+3\bar N+1
 \le11\bar N^3,
\]
and $\bar N^4\le2^{4\bar N}$. Hence
\[
 \bar N^4\bar g^4\le2^{4\bar N+44\bar N^3}(1+|\Gamma|)^4
 \le2^{48\bar N^3}(1+|\Gamma|)^4,
\]
which proves the deterministic bound of \cref{prop:summary-search}. The
nondeterministic bound is the count $\mathsf{Sp}$ above, and soundness and
completeness show that the search accepts iff $\mathcal M_0$ has a nonempty
hole of depth at most $d$.
\qed

\subsection{Accounting for the Upper Bounds}

\begin{proof}[Proof of \cref{thm:complexity}]
We account for each part of the procedure.

\emph{Satisfiability and transfers.} By \cref{lem:poly-tests}, decide in
polynomial time whether $\phi$ is satisfiable in $\mathcal D$. If it is not,
the entailment holds. Otherwise $\phi$ has a solution in arbitrary trees, and
by \eqref{eq:g7} and \eqref{eq:g8} entailment in $\mathcal D$ coincides with
entailment in arbitrary trees. By \eqref{eq:g13} and \cref{thm:holes}, the
latter holds iff neither label reader has a hole.

\emph{Reader and monitor tables.} Each label reader has $N=2^{4k+1}$ states
over $\Dir$, with $|\Dir|=n$. Its transitions, rejected states, and admitted
pairs fill $Nn(4k+1)+N+N^2$ bits, and each entry is computed in polynomial
time by \cref{lem:poly-tests}. The lettered monitor of the reader
(\cref{sec:search}) has the same transitions, the targets
$\Tfin(\zeta,a)\iff\delta(\zeta,a)\in\Rej$, and $\Lambda=\Lfin$ given by
$\Adm$ on the denoted states; it has $1+Nn$ labels. All these tables are
written in space $2^{O(k)}\operatorname{poly}(\insize\phi)$.

\emph{The empty word.} The empty word has no comparison, so it is a hole iff
the initial state is rejected; this is tested directly.

\emph{The depth bound.} As in the proof of \cref{thm:hole-decision},
\cref{thm:brt,thm:depth} show that a reader with the rejected-path property
has a hole iff the empty word is a hole or the reader has a nonempty hole of
depth at most $d=8N^2-2N-1$. Both label readers have this property
(\cref{lem:between}), and their nonempty holes are those of their lettered
monitors.

\emph{The summary search.} For each side, \cref{prop:summary-search}, with
these $N$ and $d$ and with $|\Gamma|=n$, decides the existence of such a
nonempty hole in deterministic space $2^{O((N+d+1)^3)}(1+n)^4$. Here
$N+d+1=8N^2-N\le2^{8k+5}$, so $(N+d+1)^3\le2^{24k+15}$, and the space is
$2^{2^{O(k)}}(1+n)^4$.

\emph{Total.} The stages run one after another, and the space of the whole
procedure is at most the sum of their bounds, which is
$2^{2^{O(k)}}\operatorname{poly}(\insize\phi)$. This lies in 2-EXPSPACE
because $k\le\insize\phi$. A deterministic machine that uses space
$s\ge\log\insize\phi$ has at most $2^{O(s)}$ configurations, which gives
deterministic time $2^{2^{2^{O(\insize\phi)}}}$. For the general syntax at
positive arity, \cref{prop:flattening} produces, in polynomial time, a flat
system and a variable query with the same entailment in each domain: it adds
a constant number of variables per subexpression and per encoded inequality,
and $O(n)$ symbols per encoding. The flat system has $O(\insize\phi)$
variables, so the space is $2^{2^{O(\insize\phi)}}$. At arity zero,
\cref{prop:arity-zero} gives polynomial time.
\end{proof}

\section{Constructors of Arity at Most One}
\label{app:small-arity}

\Cref{sec:complexity} concentrates on constructors of arity $n\ge2$. Its
upper bound, \cref{thm:complexity}, also covers $n\le1$, but stronger bounds
hold there: entailment is in PSPACE at arity one (\cref{thm:arity-one}) and in
polynomial time at arity zero (\cref{prop:arity-zero}), in each of the three
domains. We use the input size $\insize\phi$ and the polynomial-time tests of
\cref{lem:poly-tests}.

\subsection{Arity One}

For $n=1$ the alphabet is $\Dir=\{0\}$. Every word is a power $0^L$, every
pair of cuts $s<e\le L$ is a comparison, because $0^{L-e}\pre0^{L-s}$, and
every nonempty word is unary and so has depth zero. \Cref{thm:search} with
$d=0$ gives the length bound $\holeBound(N,0)=N+1$ of \eqref{eq:holebound}.
The following lemma improves it and needs no rejected-path property.

\begin{lemma}[One-letter readers]\label{lem:one-letter}
Let $(\mathcal M,\Adm,\Rej)$ be a reader with $N$ states over a one-letter
alphabet, and let $w$ be a hole. Then $\wpre{w}{J(w)}$ is a hole and
$J(w)\le N-1$. In particular, a reader of this kind that has a hole has a
hole of length at most $N-1$.
\end{lemma}

\begin{proof}
Let $J=J(w)$; it exists because $w$ is rejected. The prefix $\wpre{w}{J}$
ends in the rejected state $q_J$. Its comparisons are the pairs $s<e\le J$,
which are comparisons of $w$ with the same cut states; none of them is
admitted, so $\wpre wJ$ is a hole. Suppose $q_i=q_j$ for some $i<j\le J$.
With one letter, $q_{i+t}=q_{j+t}$ for every $t\ge0$. For $t=J-j$ this gives
$q_{J-(j-i)}=q_J\in\Rej$ with $J-(j-i)<J$, contradicting the minimality of
$J$. Hence $q_0,\ldots,q_J$ are pairwise distinct, and $J+1\le N$.
\end{proof}

The bound is attained among general readers: take states $0,\ldots,N-1$,
let the letter map $i$ to $\min(i+1,N-1)$, reject only $N-1$, and admit no
pair. Its only holes are the words of length at least $N-1$.

The lemma also gives an algorithm that tests a single word per reader. All
words $0^L$ follow one run $q_0,q_1,\ldots$, so every hole has the same first
rejected cut $J$, the first cut of this run with a rejected state. Follow
the run to $J$; if none of $q_0,\ldots,q_{N-1}$ is rejected, there is no
hole. Otherwise a hole exists iff $0^J$ is a hole.

\begin{theorem}[Arity one]\label{thm:arity-one}
Let $n=1$, with either variance, and let $\mathcal D$ be any of the three
domains. If $\phi$ is satisfiable in $\mathcal D$ and
$\phi\not\models_{\mathcal D,c}x\le y$, then some unsafe word has length at
most $2^{4k+1}-1$. Entailment at arity one is decidable in polynomial space
and in deterministic time $2^{O(k)}\operatorname{poly}(\insize\phi)$. In
particular, it is in PSPACE in each domain, also for the general input syntax
of \cref{prop:flattening}.
\end{theorem}

\begin{proof}
By \cref{lem:poly-tests}, satisfiability in $\mathcal D$ is decided in
polynomial time; if it fails, the entailment holds. Otherwise, by
\eqref{eq:g7}, \eqref{eq:g8}, \eqref{eq:g13}, and \cref{thm:holes}, the
entailment fails iff some label reader has a hole, and every hole is an
unsafe word. By \cref{lem:one-letter}, a label reader with a hole has one of
length at most $N-1=2^{4k+1}-1$, which proves the first claim.

For the algorithm, consider each label reader and its run
$q_0,q_1,\ldots$ on the powers of $0$. A state of \eqref{eq:g38}, two sets of
signed variables and a latch bit, takes $4k+1$ bits, and a counter of $4k+1$
bits counts up to $N-1$. Follow the run for at most $N-1$ steps, keeping only
the current state and the counter. If none of $q_0,\ldots,q_{N-1}$ is
rejected, the reader has no hole, since a hole $w$ would have
$q_{J(w)}\in\Rej$ with $J(w)\le N-1$ by \cref{lem:one-letter}. Otherwise let
$J$ be the first rejected cut of the run. Every hole $w$ is a power of $0$
with $J(w)=J$, so by \cref{lem:one-letter} the reader has a hole iff $0^J$ is
a hole, that is, iff $\neg\Adm(q_s,q_e)$ for all $s<e\le J$. The pairs $s<e$
are enumerated by two counters, and the states $q_s$ and $q_e$ are
recomputed from $q_0$. Each transition, each rejection test, and each
admission test takes polynomial time by \cref{lem:poly-tests}. Only the
closure, a constant number of states and counters, and the work space of
these tests are stored, so the space is polynomial in $\insize\phi$. The
procedure performs $O(N^3)$ transitions and tests, which gives the time
bound. For the general syntax at positive arity, \cref{prop:flattening} is a
polynomial-time reduction to a flat system with $O(\insize\phi)$ variables,
so entailment remains in PSPACE.
\end{proof}

\subsection{Arity Zero}

For $n=0$ the reader search of \cref{fig:algorithm} examines only the empty
word, and general systems cannot be reduced to our restricted flat syntax by
the positive-arity encoding of \cref{prop:flattening}. Both cases are
nevertheless easy.

\begin{proposition}[Arity zero]\label{prop:arity-zero}
Let $n=0$. The three domains coincide, and entailment is decidable in
polynomial time, both for flat systems and for general systems of
inequalities between variables and the constants $\bot,f(),\top$, with a
query between two such expressions.
\end{proposition}

\begin{proof}
At arity zero every tree is $\bot$, $f()$, or $\top$ (\cref{sec:words}), so
the three domains coincide, and by \eqref{eq:g3} the order is the chain
$\bot<f()<\top$. For a flat system, the procedure of \cref{fig:algorithm}
computes the closure, tests for a label clash, and then tests the empty
word, that is, whether the initial state of either reader is rejected. All of
this takes polynomial time by \cref{lem:poly-tests}. For a general system
$\phi$ and query $t\le t'$, the query fails under an assignment $\sigma$ iff
$\sigma(t)>\sigma(t')$. For each of the three pairs of constants $a>b$, test
whether $\phi$ together with $t\le a$, $a\le t$, $t'\le b$, and $b\le t'$ is
satisfiable; entailment holds iff all three tests fail.

Satisfiability of a finite set $\Phi$ of inequalities between variables and
constants over the chain is decided by propagation. Start with
$\lambda(v)=\bot$ for every variable, and let $\lambda$ map each constant to
itself. While some inequality $t_1\le v$ with a variable $v$ has
$\lambda(t_1)>\lambda(v)$, raise $\lambda(v)$ to $\lambda(t_1)$. Each variable
is raised at most twice. Induction on the raises shows that
$\sigma\ge\lambda$ pointwise for every solution $\sigma$. At termination,
every inequality into a variable holds under $\lambda$. If $\lambda$
satisfies $\Phi$, then $\Phi$ is satisfiable. Otherwise a violated
inequality $t_1\le b$ has a constant right side and $\lambda(t_1)>b$; every
solution $\sigma$ has $\sigma(t_1)\ge\lambda(t_1)>b$, so $\Phi$ is
unsatisfiable. All steps take polynomial time.
\end{proof}
